\section{Paired distances in actual tape cones}\label{sec:nb-section}

The recovery relation extracted from a computation is directed: an input
coordinate may occur in the cone of an output without the converse
occurrence.  The appropriate geometric object retains separate copies of
the input and output ports.  We prove that, after a paid bound on input
reuse, almost every matched pair of ports is far apart in this incidence
graph.  No assumption that the computation transposes is needed.

\subsection{The computation DAG and the incidence graph}

Fix integers $d\ge2$ and $b\ge1$, and write
\[
 L=d^2b,\qquad W=db,\qquad [L]=\{0,\ldots,L-1\}.
\]
Identify a coordinate $(i,j,s)$, with $0\le i,j<d$ and $0\le s<b$, with
the raw address $(id+j)b+s$.  The transpose permutation is
\begin{equation}\label{eq:nb-transpose}
 \Pi\bigl((id+j)b+s\bigr)=(jd+i)b+s.
\end{equation}
Throughout this section, logarithms and entropies have base two.

Consider an execution with one read-only input tape, $k\ge1$ work tapes,
and exactly $L$ output appends in $T$ steps.  Each head starts at zero and
moves by at most one cell per step.  At most one bit is appended per step.
Write $x_0(t)$ for the input-head position and $x_j(t)$ for work-head
position $j$ immediately before step $t$.  A work cell is written at every
visit, including a visit which preserves its symbol.  Let
\[
 \operatorname{prev}_j(t)=\max\{u<t:x_j(u)=x_j(t)\}
\]
when this set is nonempty.  We first take the schedule to be oblivious:
these positions, the append times, and $T$ depend on the input length
alone.  The structural argument will also apply to a single frozen
adaptive trace.

\begin{definition}[Whole-step cones]\label{def:nb-cones}
The value $v_t$ of step $t$ consists of the state after that step, all
symbols written on the $k$ work tapes, and the appended output bit or
a symbol indicating no append.  Fix a constant integer $\beta$ large
enough to encode one such value.  A genie $\Gamma\subseteq[0,T)$
reveals $(v_t)_{t\in\Gamma}$ and has cost $\beta|\Gamma|$ bits.

The computation DAG $H$ has vertices $0,\ldots,T-1$, state edges
$t\to t-1$ for $1\le t<T$, and cell edges
$t\to\operatorname{prev}_j(t)$ whenever the
displayed predecessor exists.  Parallel edges may be identified.  All
edges point backward in time.

A path $t_0\to\cdots\to t_m$ is $\Gamma$-free if
$t_1,\ldots,t_m\notin\Gamma$; its first vertex may belong to $\Gamma$.
The genie has \emph{full depth at most $D$} if every such path has at most
$D$ edges.  Let $C_\Gamma(t)$ be the set of endpoints of $\Gamma$-free
paths from $t$.  If output address $\ell$ is appended at step $t_\ell$,
put
\[
 I_\Gamma(\ell)=\{x_0(u):u\in C_\Gamma(t_\ell)\}\cap[L],
 \qquad N(q)=I_\Gamma(\Pi(q)).
\]
\end{definition}

These are the whole-step conventions of
\cite[the dependency DAG and cone lemma]{BrendeCurved}.  In particular,
a length-zero path is permitted, even when its vertex is revealed.
Here $N(q)$ is a syntactic input set; its definition does not require
that output $\Pi(q)$ equal input $q$.

Each vertex of $H$ has both in-degree and out-degree at most $k+1$.
For the in-degree assertion, a given visit has at most one next visit to
the same cell on each work tape, and at most one next state edge.
Consequently, if the genie has full depth at most $D$, then
\begin{equation}\label{eq:nb-cone-bound}
 |N(q)|\le B,\qquad B=\sum_{h=0}^{D}(k+1)^h.
\end{equation}

\begin{definition}[Paired incidence graph]\label{def:nb-incidence}
The graph $E$ has distinct vertices $L_p,R_p$ for every $p\in[L]$, with
an undirected edge $L_pR_q$ precisely when $p\in N(q)$.  The physical
port represented by $L_p$ is input address $p$; that represented by $R_q$
is output address $\Pi(q)$.  For $Q\subseteq[L]$, let $E_Q$ be the
induced graph on $\{L_p,R_p:p\in Q\}$.  Distances between vertices in
different components are infinite.
\end{definition}

The graph $E$ is an auxiliary bipartite graph on physical ports.  It is
neither the computation DAG $H$ nor the undirected graph obtained by
forgetting the directions of the recovery relation on $[L]$.  Its
bipartiteness is automatic; the assertion below concerns the distances
between the two copies of the \emph{same} coordinate.

\begin{theorem}[Paired-distance bound]\label{thm:nb-main}
\notformalized
Fix $k\ge1$, $D\ge0$, $c>0$, a degree bound $\Delta$, and an integer
$r\ge1$.  Suppose $T\le c(L+1)$ and $\Gamma$ has full depth at most
$D$.  For every $Q\subseteq[L]$ such that $E_Q$ has maximum degree at
most $\Delta$,
\begin{equation}\label{eq:nb-main-bound}
 \bigl|\{p\in Q:\operatorname{dist}_{E_Q}(L_p,R_p)\le r\}\bigr|
   =O_{k,D,c,\Delta,r}\!\left(\frac{L}{\log_2d}\right).
\end{equation}
The constant is uniform over $b\ge1$, the schedule, the genie, and $Q$.
Neither correctness as a transposer nor a bound on $|\Gamma|$ is assumed.
\end{theorem}

The proof uses a geometric relaxation in which selected paths may share
a location and subsequently separate.  These paths are auxiliary objects;
they do not assert that a coded tape symbol carries independent source
bits.  A row-entropy potential controls partial transpose transport in
this relaxation.  Alternating forward and reverse realizations of cone
matchings then converts short paired paths in $E_Q$ into that transport.

\subsection{A partial-transport inequality}

We specify the geometric relaxation so that its use does not import a
routing assumption about the original machine.  There are $K$ tapes,
one unit-speed head on each, and $\rho$ registers without spatial
addresses.  A geometric tick first permits arbitrary transfers among
the current head cells and registers, leaving other locations unchanged,
and then moves each head by at most one cell.  Each selected source
label follows a single trajectory.  Different trajectories may occupy
the same slot and later separate.  The \emph{occupancy} is the number
of selected trajectories at a slot, without regard to encoded values.
We bound occupancy at each placement before and after a local transfer.
Moving heads does not change the locations of the selected labels.

\begin{lemma}[Partial transpose transport]\label{lem:nb-transport}
\notformalized
Let $S\subseteq[L]$ have $m>0$ elements.  Suppose geometric trajectories
start at the distinct input cells $p\in S$, with no selected labels
elsewhere, and end at the distinct cells $\Pi(p)$ of one output tape.
Assume all heads start at zero, every slot has occupancy at most $C$,
and the trajectories use $U$ geometric ticks.  Then
\begin{equation}\label{eq:nb-transport-sharp}
 m\log_2d-m\log_2\!\left(\frac{L+2W}{m}\right)-m
       \le CK(K+\rho+3)U.
\end{equation}
In particular, including $m=0$,
\begin{equation}\label{eq:nb-transport-simple}
 m\log_2d\le CK(K+\rho+3)U+3L.
\end{equation}
\end{lemma}

\begin{informalproof}
Fix all trajectories before choosing any partition.  For each
$\theta\in\{0,\ldots,W-1\}$, partition every tape into the bins
$\lfloor(z-\theta)/W\rfloor$, retaining the tape name as part of the
bin name.  Merge the $K$ bins currently containing heads, together
with all registers, into one active class.  The remaining bins are
separate classes.  The active class contains at most $KW+\rho$ slots.

Choose a source label $V$ uniformly from $S$, and let
$A=\lfloor V/W\rfloor$ be its original input row.  If
$Z_{t,\theta}$ is the class containing its trajectory at boundary $t$,
define
\[
 \Phi_t=m\,\mathbb E_\theta\mathrm H(A\mid Z_{t,\theta}).
\]
This is entropy of a random \emph{source label}, not entropy of the
machine input or the contents of a cell.  For any two class variables
$Z,Z'$, the chain rule gives
\begin{equation}\label{eq:nb-chain-rule}
 \mathrm H(A\mid Z')-\mathrm H(A\mid Z)
       \le \mathrm H(Z\mid Z').
\end{equation}
Thus a new class containing $n$ labels contributes at most $n$ times
the entropy of their old class names to the increase of the unnormalized
potential.

A local transfer changes no label's class: both of its endpoints lie
in the active class.  Analyze a simultaneous movement of the heads as
$K$ successive movements, with all labels held at their physical slots.
A unit movement of a particular head crosses a bin boundary for exactly
one offset $\theta$.  At that offset, the new active class has at most
two old class names, namely the old active class and the newly joined
bin.  A released bin has only the old active name.  All other classes
are unchanged.  By occupancy, the new active class contains at most
$C(KW+\rho)\le C(K+\rho+3)W$ selected labels.  Equation
\eqref{eq:nb-chain-rule} therefore bounds the increase, after averaging
over $\theta$, by $C(K+\rho+3)$.  A pause has zero cost.  Summing gives
\begin{equation}\label{eq:nb-potential-motion}
 \Phi_U-\Phi_0\le CK(K+\rho+3)U.
\end{equation}

Initially, a shifted input bin meets at most two aligned input rows.
Only the input tape contains selected labels, and merging headed bins
merges no additional occupied input bin.  Hence $\Phi_0\le m$.

For the endpoint estimate, first keep all output bins separate.  At
output position $v=\Pi(p)$ the original source row is
$\lfloor(v\bmod W)/b\rfloor$.  This label is $W$-periodic, with exactly
$b$ occurrences of each row in a period.  Thus a shifted $W$-cell bin
containing $n$ selected labels has row multiplicities $n_a\le b$, and
\[
 n\,\mathrm H(A\mid\text{this bin})
     =n\log_2 n-\sum_a n_a\log_2n_a
     \ge n\log_2 n-n\log_2b.
\]
Use $0\log_20=0$.  At most $L/W+2$ shifted bins meet $[L]$.
Convexity of $x\log_2x$ consequently gives
\[
 m\,\mathrm H(A\mid\text{separate output bins})
   \ge m\log_2\!\left(\frac{m}{b(L/W+2)}\right)
   =m\log_2d-m\log_2\!\left(\frac{L+2W}{m}\right).
\]
Merging headed bins is a coarsening of the partition, so it cannot
decrease this conditional entropy.  This is a lower bound for $\Phi_U$.
Together with \eqref{eq:nb-potential-motion} and $\Phi_0\le m$, it proves
\eqref{eq:nb-transport-sharp}.

For $a>0$ and $0<x\le a$,
$x\log_2(a/x)\le a/(e\ln2)$.  Since $L+2W\le2L$ and $m\le L$,
\[
 m+m\log_2\!\left(\frac{L+2W}{m}\right)
 \le \left(1+\frac{2}{e\ln2}\right)L<3L.
\]
Rearranging proves \eqref{eq:nb-transport-simple}.
\end{informalproof}

The factor $b$ in the endpoint estimate is essential.  The statement
holds for arbitrary selected source subsets and every entry width; no
conversion of a word instance into a bit instance is used.

\subsection{Realizing and reversing incidence matchings}

\begin{lemma}[A matching of cone incidences]\label{lem:nb-matching}
\notformalized
Let $\mathcal M$ be a matching in $E$.  There is a geometric realization
of its edges, carrying a label from input cell $p$ to output cell
$\Pi(q)$ for each $L_pR_q\in\mathcal M$, with $K=k+2$ tapes, one
register, occupancy at most $B+1$, and $T$ ticks.  All such realizations
use the same head configurations $h_0$ at entry and $h_T$ at exit.
The geometric paths can be reversed in $2T$ ticks with the same
occupancy, from configuration $h_T$ to $h_0$.
\end{lemma}

\begin{informalproof}
For every edge $L_pR_q$ choose one $\Gamma$-free backward path from
the output step $t_{\Pi(q)}$ to a step reading raw input $p$.  It has
at most $D$ edges.  Any vertex of the computation DAG belongs to at
most $B$ chosen paths: the distinct output steps of paths using that
vertex lie within its forward radius $D$, which contains at most $B$
vertices.  This counts output endpoints, not the number of possible
combinatorial paths.

For the geometric realization, regard the input tape and the
append-output tape as two additional work-capable tapes.  The input
head follows its original trajectory.  The output head occupies the
next append address and advances exactly when an append occurs.
The additional register represents a whole step's record for purposes
of the geometric interface; it is not restricted to the original
control-state component.

Read each chosen DAG path in forward chronological order.  Its label
waits at its raw input cell until the chosen input read.  A state edge
is realized by waiting in the record register from one step to the
next.  A cell edge is realized by waiting in that work cell between
two consecutive visits.  At a path vertex, the incoming head cell or
register and the required outgoing head cell or register are present
at the local interface, so the geometric model permits the transfer.
At the chosen output step the label moves to the current output cell
and remains there.  These operations occur before the original head
movements, and therefore use the original $T$ ticks.

During a work-cell waiting interval all labels there use its same
previous-visit vertex, so there are at most $B$.  Labels in the record
register use one step vertex and satisfy the same bound.  A raw input
cell contains only its own not-yet-departed source label; a completed
output cell contains only its own output label.  The matching supplies
distinct sources and outputs, and the input and output tapes are
distinct.  The bound $B+1$ is thus conservative at every placement.
At the end of the block, all selected labels occupy output ports and
none remains in a work cell or register.

The record register is an augmentation of the geometric model.  A
syntactic state edge of the whole-step DAG may refer to components that
the real control state does not retain.  The argument does not identify
that edge with an independent information channel in the original
machine.  Instead it uses an additional register in precisely the
geometric model covered by \cref{lem:nb-transport}.

To reverse a tick, first preserve all labels and undo the original unit
head movements in one tick.  In a second tick, keep the heads fixed and
reverse the selected local transfers.  Both endpoints of each transfer
are again current head cells or registers.  This reverses the geometric
trajectories, including waiting intervals, with the same occupancy.
It does not invert a transition function or reconstruct a past machine
state.  All matchings use the same original head schedule, including
the two port tapes.  Their forward and reverse configurations therefore
agree exactly at the stated block boundaries.
\end{informalproof}

\begin{lemma}[Disjoint short paired paths]\label{lem:nb-disjoint}
\notformalized
For any fixed odd $s\ge1$, the number $m$ of vertex-disjoint $s$-edge
paths in $E$ whose endpoints are $L_p,R_p$ satisfies
\begin{equation}\label{eq:nb-disjoint-bound}
 m\le A_s\frac{L}{\log_2d},\qquad
 A_s=4sc(B+1)(k+2)(k+6)+3.
\end{equation}
\end{lemma}

\begin{informalproof}
Write $s=2a+1$ and orient each incidence path from its left endpoint
to its right endpoint.  The edges at a fixed position along these
paths form a matching, because the paths are vertex-disjoint.  Realize
the first matching forward, the second backward, and continue
alternating, ending with a forward block.

For example, an incidence path beginning
$L_p,R_v,L_u,R_w$ first carries its label from input $p$ to output
$\Pi(v)$, then from that same output cell back to input $u$, and
then to output $\Pi(w)$.  Every splice identifies the same physical
tape and address.  In particular, $\Pi$ is never performed as a free
operation; it only names which output cell a right port represents.
Forward blocks end at $h_T$, where reverse blocks begin, and reverse
blocks end at $h_0$, where the next forward blocks begin.  No head
repositioning is omitted.  At each junction the selected labels occupy
distinct ports, so occupancy does not accumulate across stages.

The resulting trajectories carry each original input $p$ to output
$\Pi(p)$, starting with every head at zero.  They use $K=k+2$ tapes,
$\rho=1$, occupancy $B+1$, and
\[
 U=(3a+1)T\le2sT\le4scL
\]
ticks.  Apply \eqref{eq:nb-transport-simple} to obtain
\eqref{eq:nb-disjoint-bound}.
\end{informalproof}

\begin{informalproof}[Proof of \cref{thm:nb-main}]
For each $p\in Q$ at paired distance at most $r$, choose one shortest
path in $E_Q$.  Its length is odd.  Put
\[
 V_r=\sum_{j=0}^{r}\max(1,\Delta)^j.
\]
A fixed vertex $z$ belongs to at most $V_r$ chosen paths: the distinct
left endpoints of paths through $z$ are within distance $r$ of $z$,
and there is one chosen path per left endpoint.

For each odd $s\le r$, choose a maximal vertex-disjoint subfamily among
the chosen paths of length $s$.  Every candidate meets a selected path.
A selected path has $s+1$ vertices and therefore meets at most
$(s+1)V_r$ candidates.  \cref{lem:nb-disjoint} gives the explicit bound
\begin{equation}\label{eq:nb-packing-bound}
 \bigl|\{p\in Q:\operatorname{dist}_{E_Q}(L_p,R_p)\le r\}\bigr|
 \le \sum_{\substack{1\le s\le r\\
s\text{ odd}}}
       (s+1)V_r A_s\frac{L}{\log_2d}.
\end{equation}
All parameters in the sum are fixed independently of $b$, the schedule,
and the genie.  This proves the theorem.
\end{informalproof}

\subsection{Input reuse, conditioning, and adaptive traces}

\begin{corollary}[Self-demands]\label{cor:nb-self-demands}
\notformalized
Without any bound on left incidence degrees,
\[
 |\{p\in[L]:p\in N(p)\}|\le A_1\frac{L}{\log_2d}.
\]
\end{corollary}

\begin{informalproof}
The edges $L_pR_p$ for self-demanding coordinates form a matching.
Apply \cref{lem:nb-disjoint} with $s=1$.
\end{informalproof}

\begin{corollary}[Paid degree reduction]\label{cor:nb-trim}
\partformal{IncidenceTrimming.lean}{paid_degree_trimming}{The native cone-size premise and the subsequent paired-distance asymptotic application.}
For an integer $M\ge1$, let
\[
 H_M=\{p\in[L]:\deg_E(L_p)>M\},\qquad Q_M=[L]\setminus H_M.
\]
Then $|H_M|\le BL/M$, and $E_{Q_M}$ has maximum degree at most
$\max(B,M)$.  For each fixed $\delta\in(0,1)$, choosing the fixed
threshold $M\ge B/\delta$ retains at least $(1-\delta)L$ coordinates
and gives \eqref{eq:nb-main-bound} at every fixed radius.
\end{corollary}

\begin{informalproof}
By \eqref{eq:nb-cone-bound},
$\sum_p\deg_E(L_p)=\sum_q|N(q)|\le BL$.
This bounds the number of left degrees exceeding $M$.  Deleting both
ports of each such coordinate cannot increase a degree.  Surviving
left degrees are at most $M$, and all right degrees were at most $B$.
\end{informalproof}

The deletion in \cref{cor:nb-trim} costs a small \emph{fixed}
fraction for a fixed threshold; it is not an uncharged $o(L)$ deletion.
Its coding interpretation is paid conditioning as in
\cref{lem:pre-conditioning}.  \cref{prop:open-transfer} below
spells out the resulting conditional transfer of an all-code half bound.

\begin{remark}[Exactly what is structural]\label{rem:nb-scope}
The premise of \cref{prop:open-transfer} is an additional
coding statement.  \cref{thm:nb-main} alone supplies no code-rate
bound.  It also supplies no genie: the depth $D$ is assumed fixed before
$d$ tends to infinity.  Its constants may grow rapidly with $D$.

Every argument proving \cref{thm:nb-main} applies to one fixed legal
adaptive execution trace with the stated time, output, and depth
properties.  Freeze that trace's head walks and append times and use
them as the geometric schedule.  The only random object in
\cref{lem:nb-transport} is an auxiliary source label, so no distribution
of inputs realizing the trace is required.  This observation does not
produce a common cone relation across inputs, bound the entropy of a
trace, or make the input coordinates independent on a trace fibre.
Those are separate obligations for an adaptive all-input coding
reduction.
\end{remark}

\subsection{Why symmetrizing the recovery relation loses the conclusion}

The distinction between incidence distance and ordinary odd cycles is
necessary even for simple native operations.  The example below forces
directed triangles in every cone relation while respecting the geometric
conditions already proved.

\begin{proposition}[Persistent directed triangles]\label{prop:nb-triangles}
\notformalized
For canonical widths $b=\max(1,\lceil\log_2d\rceil)$, there is a fixed
uniform deterministic oblivious native machine with $O(L)$ runtime
which appends exactly $L$ bits and computes
\[
 Y[j,i,s]=X[j,-j-i\bmod d,s].
\]
For every genie $\Gamma$, its demand-labelled cone relation contains
\begin{equation}\label{eq:nb-triangle-map}
 F(i,j,s)=(j,-i-j\bmod d,s)\in N(i,j,s).
\end{equation}
The map $F$ has $\gcd(3,d)b$ fixed points and otherwise consists of
disjoint directed three-cycles.  Thus the undirected symmetrization of
the recovery relation contains
$(L-\gcd(3,d)b)/3$ vertex-disjoint triangles.  Making it triangle-free
by coordinate deletion requires at least this many deletions.

Moreover, for each fixed $\varepsilon>0$, the machine admits genies of
cost at most $\varepsilon L$ and full depth bounded independently of
$d$.  The same triangle conclusion holds for those genies.
\end{proposition}

\begin{informalproof}
Output rows $j$ in order and, within each row, output entries $i$ in
order.  Obtain the entry directly from input row $j$, column
$c=-j-i\bmod d$.  Read its $b$ bits forward, appending each bit in the
same transition that reads it.  Between successive entries the target
column decreases by one.  An ordinary repositioning costs $O(b)$
input-head steps, and the single wrap within a row costs $O(db)$.
Moving from the final target of one row to the initial target of the
next costs $O(db)$.  The total input motion and emission cost is
therefore $O(d^2b)$.

Here is a fixed-resource implementation of the indexing and motion
control.  First count the raw length $L$ with a little-endian binary
counter, returning its head to zero after each increment.  The total
carry distance is $O(L)$.  Rewind the input head, and recover $d,b$
from the strictly increasing function
$d\mapsto d^2\max(1,\lceil\log_2d\rceil)$ by binary search and
ordinary arithmetic on $O(\log L)$ bits.  This setup takes
$\operatorname{poly}(\log L)=O(L+1)$ time on a fixed number of tapes.
Noncanonical lengths may use an immediate fallback after this test.

At canonical lengths $d\ge2$, every subsequent index and displacement
has $O(b)$ bits.  Per-entry operations are addition, subtraction,
comparison, and modular decrement, implementable in $O(b)$ time by
full scans of a fixed number of binary registers.  A reusable unary
guide of length $b$ controls blocks of $b$ input moves or output
appends.  Each use scans outward and returns, with the input head
paused during the return.  Longer displacements use $O(d)$ such
blocks with $O(b)$ counter work per block; they cost $O(db)$ each.
Hence all entry operations, row wraps, and row transitions still cost
$O(L)$.  The work tapes are registers or guides accessed by monotone
outward-and-return scans from zero, so their walks are home-based.
All motion depends on the public length and counters; payload bits
affect only the emitted values.  This proves uniformity, obliviousness,
and the claimed native cost with a fixed tape alphabet and control.

The step emitting physical output $(j,i,s)$ reads raw input
$(j,-j-i\bmod d,s)$.  The length-zero path at this step is
$\Gamma$-free for every $\Gamma$, including one revealing that step.
Demand relabelling by \eqref{eq:nb-transpose} proves
\eqref{eq:nb-triangle-map}.  Direct calculation gives
\[
 F^2(i,j,s)=(-i-j,i,s),\qquad F^3(i,j,s)=(i,j,s).
\]
Its fixed points satisfy $i=j$ and $3i=0$ in $\mathbb Z/d\mathbb Z$,
so there are exactly $\gcd(3,d)b$ of them.  All other orbits have size
three and give vertex-disjoint triangles after symmetrization.

Finally use the home-based depth-reduction theorem and long-edge
conversion of \cite[Theorem~B and the proof of Corollary~B]{BrendeCurved}.
For fixed machine resources and each $\varepsilon>0$, these give a
genie with cost at most $\varepsilon L$ and constant full depth.
That construction uses the schedule and its runtime, not transpose
correctness.  It therefore applies to this machine.  The triangles
persist because \eqref{eq:nb-triangle-map} held for every genie.
\end{informalproof}

The guaranteed arcs in this example contribute only the incidence
matching $L_{F(q)}R_q$; they do not by themselves connect $L_q$ to
$R_q$ around the directed three-cycle.  There is therefore no conflict
with \cref{thm:nb-main}.  Additional actual cone edges are permitted
throughout the argument.  The example rules out a universal local
bipartiteness conclusion for the symmetrized recovery graph, while
leaving the directed paired-distance condition intact.
